\worksheettitle
  {Ломаная и многоугольник}
  {\modulemark{M05} Версия учителя}

\begin{teacherbox}[Паспорт модуля]
  \textbf{Результат:} ученик различает ломаную и многоугольник, определяет
  замкнутость, называет вершины, звенья и стороны, обходит границу
  последовательно.

  \textbf{Граница содержания:} различение стороны и диагонали подробно
  вводится в M06.
\end{teacherbox}

\begin{teacherbox}[Методический акцент]
  Просите ученика показывать каждый названный элемент на рисунке. Число букв
  в имени незамкнутой ломаной не равно числу её звеньев: считать нужно
  реальные переходы между соседними вершинами.
\end{teacherbox}

\section{Вход и разобранный пример}

Во входе: граница замкнута; $3$ вершины; треугольник.

\openbrokenfigure

В пропусках: вершины; звенья; $5$ вершин; $4$ звена.

\section{Ответы к практике}

\studentrecognizefigure

\begin{practicebox}[1. Тип ломаной]
  a) рисунки 1 и 3; б) рисунки 2 и 4; в) $4$ вершины и $3$ звена.
\end{practicebox}

\polygonpartsfigure

\begin{practicebox}[2. Элементы шестиугольника]
  Вершины: $A,B,C,D,E,F$. Стороны: $AB,BC,CD,DE,EF,FA$.
  Шесть сторон и шесть вершин.
\end{practicebox}

\begin{practicebox}[3. Самостоятельно]
  Четырёхугольник $KLMN$; стороны $KL,LM,MN,NK$.
  Отрезок $KM$ соединяет несоседние вершины и не входит в указанную границу.
\end{practicebox}

\begin{teacherbox}[К построению]
  В задании 4 допустима любая простая последовательность $P-Q-R-S$.
  Замыкающее звено $SP$ превращает её в границу четырёхугольника.
\end{teacherbox}

\section{Ошибка и контроль}

\begin{warningbox}[Ожидаемое исправление]
  Причина: учтена только замкнутость, но не устройство границы. Окружность
  не состоит из отрезков, поэтому не является ломаной и не образует
  многоугольник.
\end{warningbox}

\openclosedfigure

\begin{checkpointbox}[Ответы]
  \begin{enumerate}[label=\textbf{\arabic*.}]
    \item Замкнута ломаная $KLMNP$ справа.
    \item Вершины слева: $A,B,C,D$; звенья $AB,BC,CD$.
    \item Справа $5$ сторон и $5$ вершин.
    \item Например, $KLMNP$; допустимы циклический сдвиг и обратный обход.
  \end{enumerate}
\end{checkpointbox}

\begin{teacherbox}[Решение о маршруте]
  M05-R назначается, если ученик считает точки вместо звеньев, не видит
  разрыва границы или называет внутренний отрезок стороной. M05-H достаточно
  при верной классификации и единичной ошибке в обозначении. Устойчивый
  результат открывает M06.
\end{teacherbox}
